% ============================================================
%  يُترجَم بـ XeLaTeX  (في Overleaf: Menu → Compiler → XeLaTeX)
%  الخط: Amiri (متوفر في Overleaf، أو حمّله من github.com/aliftype/amiri)
% ============================================================
\documentclass[12pt,a4paper]{article}
\usepackage[a4paper,margin=1.7cm,top=2.3cm,bottom=2.9cm,headheight=15pt,footskip=1.5cm]{geometry}
\usepackage{amsmath,amssymb}
\usepackage{icomma}
\usepackage[english,bidi=default,layout=lists]{babel}
\babelprovide[import,main]{arabic}
\babelfont{rm}[Scale=1.08]{Amiri}
\babelfont{sf}[Scale=1.08]{Amiri}
\usepackage[most]{tcolorbox}
\usepackage{tikz}
\usetikzlibrary{arrows.meta,calc,shadings}
\usepackage{array}
\usepackage{colortbl}
\usepackage{fancyhdr}
\usepackage{titlesec}
\usepackage{enumitem}
\usepackage{collcell}
\usepackage{needspace}
\usepackage{etoolbox}
% أعمدة جداول تُكتب خلاياها بالعربية في الاتجاه الصحيح
\newcolumntype{C}{>{\collectcell\ar}c<{\endcollectcell}}
% تجنب عنوان يتيم في أسفل الصفحة
\pretocmd{\section}{\Needspace{15\baselineskip}}{}{}
\pretocmd{\subsection}{\Needspace{11\baselineskip}}{}{}

% ---------- الألوان ----------
\definecolor{cDef}{HTML}{1F4E9C}
\definecolor{cThm}{HTML}{B3261E}
\definecolor{cRem}{HTML}{2E7D32}
\definecolor{cEx}{HTML}{C76B00}
\definecolor{cMeth}{HTML}{6A1B9A}
\definecolor{cSol}{HTML}{00695C}
\definecolor{cSec}{HTML}{1F3A68}

% ---------- العناوين ----------
\titleformat{\section}{\Large\bfseries\color{cSec}}{\thesection.}{0.6em}{}[{\color{cSec!60}\titlerule[1pt]}]
\titleformat{\subsection}{\large\bfseries\color{cThm!85!black}}{\thesubsection.}{0.5em}{}
\titlespacing*{\section}{0pt}{14pt}{8pt}
\titlespacing*{\subsection}{0pt}{10pt}{5pt}
\setlength{\parindent}{0pt}
\setlength{\parskip}{4pt}
\setlist{itemsep=2pt,topsep=3pt}

% ---------- الصناديق (تُرسم في سياق LTR ويُكتب محتواها بالعربية) ----------
\newtcolorbox{arbox}[3][]{enhanced,
  colback=#2!5,colframe=#2!80!black,arc=2mm,boxrule=0.8pt,
  left=6pt,right=6pt,top=8pt,bottom=5pt,
  fonttitle=\bfseries,coltitle=white,
  title={\foreignlanguage{arabic}{#3}},
  attach boxed title to top right={xshift=-5mm,yshift=-3mm},
  boxed title style={colback=#2!80!black,arc=1.5mm,boxrule=0pt},
  before upper={\selectlanguage{arabic}},
  before skip=8pt,after skip=8pt,breakable,#1}
\newenvironment{BOX}[2]{\begin{otherlanguage}{english}\begin{arbox}{#1}{#2}}{\end{arbox}\end{otherlanguage}}
\newenvironment{tarif}[1][تعريف]{\begin{BOX}{cDef}{#1}}{\end{BOX}}
\newenvironment{khassia}[1][خاصية]{\begin{BOX}{cThm}{#1}}{\end{BOX}}
\newenvironment{molahada}[1][ملاحظة]{\begin{BOX}{cRem}{#1}}{\end{BOX}}
\newenvironment{mithal}[1][مثال]{\begin{BOX}{cEx}{#1}}{\end{BOX}}
\newenvironment{tariqa}[1][طريقة]{\begin{BOX}{cMeth}{#1}}{\end{BOX}}

% ---------- التمارين والحلول ----------
\newcounter{exo}
\newenvironment{tamrin}[1]{\stepcounter{exo}\begin{BOX}{cSec}{التمرين \theexo: #1}}{\end{BOX}}
\newenvironment{hal}{\Needspace{12\baselineskip}\begin{otherlanguage}{english}\begin{arbox}[breakable]{cSol}{حل التمرين \theexo}}{\end{arbox}\end{otherlanguage}}

% رسومات TikZ تُوضع في سياق LTR
\newenvironment{LTR}{\begin{otherlanguage}{english}}{\end{otherlanguage}}
\newcommand{\ar}[1]{\foreignlanguage{arabic}{#1}}

% ---------- رموز ----------
% نقاط الحذف داخل الرياضيات (\dots يسبب خللاً في الاتجاه مع babel)
\renewcommand{\dots}{\mathinner{\ldotp\ldotp\ldotp}}
\renewcommand{\ldots}{\mathinner{\ldotp\ldotp\ldotp}}
\renewcommand{\cdots}{\mathinner{\cdotp\cdotp\cdotp}}
\newcommand{\N}{\mathbb{N}}
\newcommand{\Z}{\mathbb{Z}}
\newcommand{\D}{\mathbb{D}}
\newcommand{\Q}{\mathbb{Q}}
\newcommand{\R}{\mathbb{R}}
\newcommand{\pgcd}{\operatorname{PGCD}}
\newcommand{\ppcm}{\operatorname{PPCM}}

% ---------- تمثيل مجال على مستقيم عددي ----------
% \NLseg{a}{b}{[}{]}  مجال محدود
\newcommand{\NLseg}[4]{\begin{LTR}\begin{tikzpicture}[baseline=-0.6ex]
 \draw[-{Stealth[length=2mm]},gray!70](0,0)--(3.4,0);
 \draw[line width=2pt,cThm!80!black](0.7,0)--(2.7,0);
 \node[cThm!80!black,font=\LARGE] at (0.7,0){$#3$};
 \node[cThm!80!black,font=\LARGE] at (2.7,0){$#4$};
 \node[below=3pt,font=\footnotesize] at (0.7,0){$#1$};
 \node[below=3pt,font=\footnotesize] at (2.7,0){$#2$};
\end{tikzpicture}\end{LTR}}
% \NLright{a}{[}  مجال نحو +∞
\newcommand{\NLright}[2]{\begin{LTR}\begin{tikzpicture}[baseline=-0.6ex]
 \draw[-{Stealth[length=2mm]},gray!70](0,0)--(3.4,0);
 \draw[line width=2pt,cThm!80!black](1.3,0)--(3.3,0);
 \node[cThm!80!black,font=\LARGE] at (1.3,0){$#2$};
 \node[below=3pt,font=\footnotesize] at (1.3,0){$#1$};
\end{tikzpicture}\end{LTR}}
% \NLleft{b}{]}  مجال نحو −∞
\newcommand{\NLleft}[2]{\begin{LTR}\begin{tikzpicture}[baseline=-0.6ex]
 \draw[-{Stealth[length=2mm]},gray!70](0,0)--(3.4,0);
 \draw[line width=2pt,cThm!80!black](0.05,0)--(2.1,0);
 \node[cThm!80!black,font=\LARGE] at (2.1,0){$#2$};
 \node[below=3pt,font=\footnotesize] at (2.1,0){$#1$};
\end{tikzpicture}\end{LTR}}

% ---------- أيقونات مواقع التواصل (مرسومة بـ TikZ، لا تحتاج أي حزمة) ----------
\definecolor{cFB}{HTML}{1877F2}
\definecolor{cTG}{HTML}{229ED9}
\definecolor{cIG1}{HTML}{FEDA75}
\definecolor{cIG2}{HTML}{D62976}
\definecolor{cIG3}{HTML}{4F5BD5}
\newcommand{\fbicon}{\begin{LTR}\begin{tikzpicture}[baseline=-0.75ex,x=1em,y=1em]
 \fill[cFB] (0,0) circle (0.55);
 \node[white,inner sep=0,font=\usefont{TU}{lmss}{bx}{n}\small] at (0.04,-0.02) {f};
\end{tikzpicture}\end{LTR}}
\newcommand{\igicon}{\begin{LTR}\begin{tikzpicture}[baseline=-0.75ex,x=1em,y=1em]
 \shade[lower left=cIG1,lower right=cIG2,upper left=cIG2,upper right=cIG3,rounded corners=0.28em] (-0.55,-0.55) rectangle (0.55,0.55);
 \draw[white,line width=0.09em,rounded corners=0.18em] (-0.36,-0.36) rectangle (0.36,0.36);
 \draw[white,line width=0.09em] (0,0) circle (0.17);
 \fill[white] (0.22,0.22) circle (0.05);
\end{tikzpicture}\end{LTR}}
\newcommand{\tgicon}{\begin{LTR}\begin{tikzpicture}[baseline=-0.75ex,x=1em,y=1em]
 \fill[cTG] (0,0) circle (0.55);
 \fill[white] (-0.33,0.0) -- (0.32,0.27) -- (0.2,-0.3) -- (0.0,-0.15) -- (-0.08,-0.28) -- (-0.08,-0.08) -- cycle;
 \fill[cTG!70!black] (-0.08,-0.08) -- (0.25,0.2) -- (0.0,-0.15) -- (-0.08,-0.28) -- cycle;
\end{tikzpicture}\end{LTR}}
% سطر مواقع التواصل أسفل كل صفحة
% (يُكتب السطر بترتيب بصري من اليسار إلى اليمين: الأيقونة يمين كل اسم)
\newcommand{\PageBox}{\begin{LTR}\tikz[baseline=(p.base)]\node[fill=black,text=white,inner sep=2.5pt,minimum width=1.6em,minimum height=1.3em,font=\usefont{TU}{lmss}{bx}{n}\small](p){\ifnum\value{page}<10 0\fi\arabic{page}};\end{LTR}}
\newcommand{\SocialRow}{\begin{otherlanguage}{english}%
 \ar{\textbf{الأستاذ سوايسية هشام للرياضيات}}\ \tgicon\hfill
 \textbf{prof\_hicham\_math}\ \igicon\hfill
 \ar{\textbf{الأستاذ سوايسية هشام للرياضيات}}\ \fbicon\hspace{2.2em}\PageBox\end{otherlanguage}}

% ---------- رأس وتذييل الصفحة ----------
\pagestyle{fancy}
\fancyhf{}
\renewcommand{\headrulewidth}{0.6pt}
\renewcommand{\footrulewidth}{0.4pt}
% \SetHeader{يمين الرأس}{يسار الرأس}{يمين التذييل}{يسار التذييل}
\newcommand{\SetHeader}[2]{%
 \fancyhead[C]{\parbox{\textwidth}{\ar{#1\hfill #2}}}%
 \fancyfoot[C]{\parbox{\textwidth}{\small\SocialRow}}}

\pretocmd{\section}{\Needspace{24\baselineskip}}{}{}

% معلم بشبكة: \Grid{xmin}{xmax}{ymin}{ymax}
\newcommand{\Grid}[4]{%
 \draw[gray!25] (#1,#3) grid (#2,#4);
 \draw[-{Stealth}] (#1-0.2,0)--(#2+0.3,0) node[below]{$x$};
 \draw[-{Stealth}] (0,#3-0.2)--(0,#4+0.3) node[left]{$y$};
 \node[below left,font=\scriptsize] at (0,0){$O$};}

\SetHeader{\small\bfseries\color{cSec}الوحدة الأولى: الدوال العددية -- سلسلة تمارين}{\small\color{cSec}السنة الثانية ثانوي -- الشعب العلمية}

\begin{document}

\begin{LTR}
\begin{tcolorbox}[enhanced,colback=cSec,colframe=cSec,arc=3mm,halign=center,
  boxrule=0pt,top=10pt,bottom=10pt,before upper={\selectlanguage{arabic}}]
{\color{white}\Huge\bfseries سلسلة تمارين (بدون حلول)}\\[6pt]
{\color{white}\large الوحدة الأولى: الدوال العددية -- السنة الثانية ثانوي}\\[4pt]
{\color{white!85!cSec}\normalsize إعداد: الأستاذ سوايسية هشام}
\end{tcolorbox}
\end{LTR}

\begin{molahada}[توجيهات]
التمارين مرتبة حسب محاور الوحدة، ثم تمرينان شاملان في الأخير. حاول حل كل تمرين بنفسك قبل الاطلاع على حله.
\end{molahada}

% =================================================================
\section{مجموعة التعريف والعمليات على الدوال}
% =================================================================
\begin{tamrin}{مجموعة التعريف}
عيّن مجموعة تعريف كل دالة من الدوال التالية:
\[f_1(x)=\frac{2x+1}{x^2-4}\qquad f_2(x)=\sqrt{3-2x}\qquad f_3(x)=\frac{1}{\sqrt{x+1}}\]
\[f_4(x)=\frac{\sqrt{x-1}}{x-3}\qquad\qquad f_5(x)=\frac{x}{x^2+1}\]
\end{tamrin}

\begin{tamrin}{تساوي دالتين}
في كل حالة، هل الدالتان $f$ و $g$ متساويتان؟ علّل.
\begin{enumerate}
\item $f(x)=\left(\sqrt{x+2}\right)^2$ \quad و \quad $g(x)=x+2$.
\item $f(x)=\dfrac{1}{\sqrt x-1}$ \quad و \quad $g(x)=\dfrac{\sqrt x+1}{x-1}$.
\end{enumerate}
\end{tamrin}

\begin{tamrin}{العمليات على الدوال}
$f$ و $g$ دالتان معرفتان على $\R$ بـ: \ $f(x)=x^2+1$ \ و \ $g(x)=x+2$.\\
عرّف كلاً من الدوال: \ $f+3$ ، \ $-2f$ ، \ $f+g$ ، \ $f\times g$ ، \ $\dfrac fg$.
\end{tamrin}

% =================================================================
\section{تركيب الدوال واتجاه التغير}
% =================================================================
\begin{tamrin}{تركيب دالتين}
$f(x)=x^2+1$ معرفة على $\R$ ، و $g(x)=\dfrac{1}{x-2}$ معرفة على $\R-\{2\}$.
\begin{enumerate}
\item اكتب كلاً من $f$ و $g$ على شكل مركب دالتين مرجعيتين.
\item عيّن مجموعة تعريف $g\circ f$ ثم احسب $(g\circ f)(x)$.
\item عيّن مجموعة تعريف $f\circ g$ ثم احسب $(f\circ g)(x)$. ماذا تلاحظ؟
\end{enumerate}
\end{tamrin}

\begin{tamrin}{تفكيك دالة واتجاه تغيرها}
\begin{enumerate}
\item فكك كل دالة إلى مركب دالتين بسيطتين: \ $f(x)=(x+2)^2+1$ \ ، \ $g(x)=\dfrac{3}{x+1}$ \ ، \ $h(x)=\left|\dfrac{2x}{5}-1\right|$.
\item $k$ الدالة المعرفة على $]-\infty;3]$ بـ: $k(x)=\sqrt{3-x}$. اكتب $k$ على شكل مركب دالتين ثم استنتج اتجاه تغيرها.
\end{enumerate}
\end{tamrin}

\begin{tamrin}{اتجاه التغير باستعمال الدوال المرجعية}
ادرس اتجاه تغير كل دالة على المجال المعطى:
\begin{enumerate}
\item $f(x)=2\sqrt x+3$ على $[0;+\infty[$.
\item $g(x)=\dfrac1x-4$ على $]-\infty;0[$.
\item $h(x)=-\dfrac{3}{x-1}+2$ على $]1;+\infty[$.
\item $k(x)=(x+1)^2$ على $]-1;+\infty[$.
\end{enumerate}
\end{tamrin}

\begin{tamrin}{اتجاه تغير $\sqrt u$ و $\dfrac1u$}
ادرس اتجاه تغير كل دالة على المجال المعطى:
\begin{enumerate}
\item $f(x)=\sqrt{3-x}$ على $]-\infty;3]$.
\item $g(x)=\dfrac{1}{x^2-4x+5}$ على $\R$ ، ثم استنتج القيمة الحدية العظمى لـ $g$.
\item $h(x)=\sqrt{x^2+2x+3}$ على $\R$.
\item $k(x)=3-\dfrac{2}{\sqrt x+1}$ على $[0;+\infty[$.
\end{enumerate}
\end{tamrin}

% =================================================================
\section{التمثيلات البيانية}
% =================================================================
\begin{tamrin}{تمثيلات مستنتجة}
$f$ الدالة المعرفة على $\R$ بـ: $f(x)=x^2-2$.
\begin{enumerate}
\item ارسم $(C_f)$ انطلاقاً من التمثيل البياني للدالة مربع.
\item استنتج التمثيل البياني للدالة $g(x)=|f(x)|$ ، ثم للدالة $h(x)=f\big(|x|\big)$. ماذا تلاحظ؟
\end{enumerate}
\end{tamrin}

\begin{tamrin}{منحنى مستنتج بانسحاب}
$f$ الدالة المعرفة على $\R-\{3\}$ بـ: $f(x)=\dfrac{1}{x-3}+2$.
\begin{enumerate}
\item بيّن أن $(C_f)$ صورة منحنى الدالة مقلوب بانسحاب يُطلب تعيين شعاعه.
\item استنتج اتجاه تغير $f$ ، ثم ارسم $(C_f)$.
\end{enumerate}
\end{tamrin}

% =================================================================
\section{الشفعية والتناظر وتغيير المعلم}
% =================================================================
\begin{tamrin}{الشفعية}
ادرس شفعية كل دالة: \quad $f(x)=x^3-3x$ \ ، \ $g(x)=x^2+|x|$ \ ، \ $h(x)=x^2+x$.
\end{tamrin}

\begin{tamrin}{محور التناظر ومركز التناظر}
\begin{enumerate}
\item $f(x)=x^2+2x+1$. بيّن أن المستقيم ذا المعادلة $x=-1$ محور تناظر لـ $(C_f)$.
\item $g(x)=\dfrac{3x-5}{x-2}$ على $\R-\{2\}$. بيّن أن النقطة $\Omega(2;3)$ مركز تناظر لـ $(C_g)$ بطريقتين.
\end{enumerate}
\end{tamrin}

\begin{tamrin}{تغيير المعلم}
$f(x)=x^2-4x+9$ و $\Omega(2;5)$. اكتب معادلة $(C_f)$ في المعلم $\left(\Omega;\vec i,\vec j\right)$ ، ثم استنتج طبيعة $(C_f)$ ومحور تناظره.
\end{tamrin}

% =================================================================
\section{كثيرات الحدود والدرجة الثانية}
% =================================================================
\begin{tamrin}{تحليل كثير حدود}
$f(x)=x^3-x^2-4x+4$.
\begin{enumerate}
\item تحقق أن $2$ جذر لـ $f$.
\item عيّن الأعداد $a$ ، $b$ ، $c$ بحيث $f(x)=(x-2)\left(ax^2+bx+c\right)$ ، ثم حلّل $f(x)$ إلى جداء عوامل من الدرجة الأولى.
\item حل في $\R$ المعادلة $f(x)=0$ ثم المتراجحة $f(x)\ge0$.
\end{enumerate}
\end{tamrin}

\begin{tamrin}{المعادلات والمتراجحات من الدرجة الثانية}
\begin{enumerate}
\item $f(x)=x^2+3x-10$. أعط الشكل النموذجي لـ $f(x)$ ، ثم بيّن أن $f(x)\ge-\dfrac{49}{4}$ لكل $x$ ، ثم حلّل $f(x)$.
\item حل في $\R$ المعادلات: \ (أ) $2x^2+5x-3=0$ \ (ب) $-x^2+2x-5=0$ \ (ج) $x^2-8x+16=0$.
\item حل في $\R$ المتراجحات: \ (أ) $-3x^2+4x-1>0$ \ (ب) $2x^2-2x-4\le0$ \ (ج) $x^2-2\sqrt3\,x+3\le0$.
\end{enumerate}
\end{tamrin}

\begin{tamrin}{مضاعفة التربيع والمعادلات الصماء}
\begin{enumerate}
\item حل في $\R$: \ (أ) $x^4+x^2-6=0$ \ (ب) $x^4-5x^2+4=0$ \ (ج) $2x^4+5x^2+2=0$.
\item حل في $\R$ المتراجحة: $x^4-7x^2+12<0$.
\item حل في $\R$ المعادلة: $\sqrt{x+3}=x+1$.
\end{enumerate}
\end{tamrin}

\begin{tamrin}{المجموع والجداء}
\begin{enumerate}
\item أوجد عددين مجموعهما $7$ وجداؤهما $10$.
\item دون حل المعادلة $2x^2-7x+3=0$ ، بيّن أنها تقبل حلين موجبين، ثم تحقق أن $3$ حل لها واستنتج الحل الآخر.
\end{enumerate}
\end{tamrin}

% =================================================================
\section{كثيرات الحدود: القسمة الإقليدية وجدول هورنر}
% =================================================================
\begin{tamrin}{القسمة الإقليدية لكثيرات الحدود}
$P(x)=2x^3-5x^2+x+4$.
\begin{enumerate}
\item أنجز القسمة الإقليدية لـ $P(x)$ على $(x-2)$ ، ثم اكتب $P(x)$ على الشكل $P(x)=(x-2)\,q(x)+r$.
\item احسب $P(2)$ وقارنه مع الباقي $r$. \ هل العدد $2$ جذر لـ $P$؟
\item أنجز القسمة الإقليدية لـ $P(x)$ على $d(x)=x^2-x+1$.
\item استنتج من السؤال 1 حلول المعادلة $P(x)=2$ في $\R$.
\end{enumerate}
\end{tamrin}

\begin{tamrin}{جدول هورنر والتحليل}
$h(x)=x^3+2x^2-5x-6$.
\begin{enumerate}
\item تحقق أن $-1$ جذر لـ $h$.
\item باستعمال جدول هورنر ، عيّن كثير الحدود $g$ بحيث $h(x)=(x+1)\,g(x)$.
\item حلّل $h(x)$ إلى جداء عوامل من الدرجة الأولى.
\item حل في $\R$ المعادلة $h(x)=0$ ثم المتراجحة $h(x)>0$.
\end{enumerate}
\end{tamrin}

\begin{tamrin}{تعيين وسيط حتى يكون عدد جذراً}
\begin{enumerate}
\item $m$ عدد حقيقي و $P(x)=x^3+mx^2-4x+12$.
\begin{itemize}
\item[(أ)] عيّن $m$ حتى يكون $2$ جذراً لـ $P$.
\item[(ب)] من أجل هذه القيمة ، حلّل $P(x)$ باستعمال جدول هورنر.
\end{itemize}
\item $a$ و $b$ عددان حقيقيان و $Q(x)=x^3+ax^2+bx-6$.
\begin{itemize}
\item[(أ)] عيّن $a$ و $b$ حتى يكون العددان $1$ و $-2$ جذرين لـ $Q$.
\item[(ب)] حلّل عندئذ $Q(x)$ إلى جداء عوامل من الدرجة الأولى.
\end{itemize}
\end{enumerate}
\end{tamrin}

\begin{tamrin}{تحليل كامل ، إشارة ومتراجحات}
$f(x)=2x^3+x^2-13x+6$.
\begin{enumerate}
\item احسب $f(2)$ ، ثم عيّن كثير الحدود $g$ بحيث $f(x)=(x-2)\,g(x)$.
\item حلّل $f(x)$ إلى جداء عوامل من الدرجة الأولى.
\item ادرس إشارة $f(x)$ ، ثم حل في $\R$ المتراجحة $f(x)\le0$.
\item حل في $\R$ المتراجحة: $f(x)\ge(x-2)(x+3)$.
\end{enumerate}
\end{tamrin}

% =================================================================
\section{تمارين شاملة}
% =================================================================
\begin{tamrin}{شامل (1)}
$f$ الدالة المعرفة على $\R$ بـ: $f(x)=x^2-2x-3$ ، و $(C_f)$ تمثيلها البياني.
\begin{enumerate}
\item اكتب $f(x)$ على الشكل النموذجي، ثم استنتج أن $(C_f)$ صورة منحنى الدالة مربع بانسحاب يُطلب تعيين شعاعه.
\item استنتج اتجاه تغير $f$ ومحور تناظر $(C_f)$.
\item عيّن نقط تقاطع $(C_f)$ مع المحورين، ثم ادرس إشارة $f(x)$.
\item ناقش بيانياً، حسب قيم الوسيط الحقيقي $m$ ، عدد حلول المعادلة $f(x)=m$.
\item $k(x)=\sqrt{f(x)}$. عيّن $D_k$ ثم اتجاه تغير $k$.
\end{enumerate}
\end{tamrin}

\begin{tamrin}{شامل (2)}
$f$ الدالة المعرفة على $\R-\{1\}$ بـ: $f(x)=\dfrac{2x+1}{x-1}$.
\begin{enumerate}
\item بيّن أنه من أجل كل $x\ne1$: \ $f(x)=2+\dfrac{3}{x-1}$.
\item استنتج اتجاه تغير $f$ على كل من المجالين $]-\infty;1[$ و $]1;+\infty[$.
\item بيّن أن النقطة $\Omega(1;2)$ مركز تناظر لـ $(C_f)$.
\item حل في $\R-\{1\}$ المعادلة $f(x)=x$ ، وفسّر النتيجة بيانياً.
\end{enumerate}
\end{tamrin}

\vfill
\begin{center}
{\color{cSec}\large\bfseries بالتوفيق والنجاح}
\end{center}
\end{document}
